(1) is Lemma \claimref{one-relation.vector-derivative-covariant-first-order-operator}. (2) The Klein-Gordon square: $\nabla(\nabla\psi\,I\sigma_{3}) = \nabla(m\psi\gamma_{0})$, and since $I\sigma_{3}$ and $\gamma_{0}$ are constant and $\nabla$ acts on the left, $\nabla^{2}\psi\, I\sigma_{3} = m\,\nabla\psi\,\gamma_{0} = m\,(m\psi\gamma_{0}(I\sigma_{3})^{-1})\gamma_{0} = -m^{2}\psi\,\gamma_{0}\,I\sigma_{3}\,\gamma_{0} = -m^{2}\psi\, I\sigma_{3}$, using $(I\sigma_{3})^{-1} = -I\sigma_{3}$ and $\gamma_{0}\,I\sigma_{3}\,\gamma_{0} = I\sigma_{3}$ (the spin plane commutes with $\gamma_{0}$); right-multiplying by $(I\sigma_{3})^{-1}$ gives $\nabla^{2}\psi = -m^{2}\psi$, and $\nabla^{2} = \eta^{\mu\nu}\partial_{\mu}\partial_{\nu} = \Box$ because $\{\gamma^{\mu}, \gamma^{\nu}\} = 2\eta^{\mu\nu}$ (\claimref{std-algebra.vector-squares}). The script checks the three algebraic facts used. The identification of $m$ with the dictionary's mass is then the statement that the same Klein-Gordon operator appears, which is what the dictionary (\claimref{one-relation.imported-condensate-derived-content}) computes for a density of the relevant weight; the proposition claims it at that level and no higher. (2), form: Lemmas \claimref{one-relation.form-law-forced-square} and \claimref{one-relation.odd-square-roots-chiral-phase}. (3) The script checks that $(\gamma_{0}\gamma_{1})^{2} = +1$, that $(\gamma_{2}\gamma_{1})^{2} = -1$ with $\gamma_{2}\gamma_{1}$ commuting with $\gamma_{0}$ and anticommuting with $\gamma_{1}, \gamma_{2}$, and that a rational rest-space rotation carries $\gamma_{2}\gamma_{1}$ to a distinct bivector squaring to $-1$ and commuting with $\gamma_{0}$. The orientation identification is Fact \claimref{one-relation.imported-flag-arrow-superselection}, cited at its own level, under (S, \claimref{standard-model.hyp-s-slot-statistics}); the unsourced-frame clause is the observation that the de Sitter frozen point is invariant under $SO(1,4)$, which contains no preferred timelike direction, together with the absence in the series of any theorem selecting one, sharpened by Lemma \claimref{one-relation.one-datum-not-two-imaginary-unit} to the statement that none can exist without breaking the invariance; the one-datum and one-bit clauses are Lemma \claimref{one-relation.one-datum-not-two-imaginary-unit} and Corollary \claimref{one-relation.arrow-frequency-are-one-bit-level}. (4) is the standard nonrelativistic descent, imported; the script verifies the Pauli identity on nine rational vector pairs in $\Cl(3,0)$ with $I_{3}^{2} = -1$. The proposition's level is declared in its statement.
